\paragraph*{Syntax at a glance}

\begin{table}[htbp]
\centering\small
\begin{tabularx}{\linewidth}{@{}lX@{}}
\toprule
\textbf{Syntax} & \textbf{Meaning} \\
\midrule
\texttt{dma.\allowbreak{}copy.\allowbreak{}swiz\textless{}N\textgreater{} src =\textgreater{} dst} & Copy with swizzle width \\texttt{N} in bytes \\
\texttt{mma.\allowbreak{}load.\allowbreak{}swiz\textless{}N\textgreater{} src} & MMA operand load consistent with swizzle \texttt{N} \\
\texttt{.\allowbreak{}view(M, N).\allowbreak{}from(r, c)} & Arbitrary-offset \texttt{M x N} window \\
\texttt{.\allowbreak{}subspan(M, K).\allowbreak{}step(sM, sK).\allowbreak{}at(i, j)} & Strided tile selection \\
\texttt{.zfill} & Zero-fill out-of-bounds elements on copy source \\
\texttt{span\_as(dims)} & Reinterpret linear storage as shaped tensor \\
\bottomrule
\end{tabularx}
\end{table}

Section~\ref{sec:3-6} takes the tile expressions above and uses them as the source and destination of movement statements. For the cases where a program needs plain C++, Section~\ref{sec:2-13} collects the \texttt{\_\_device\_\_} functions, preprocessor guards and injected source text.
